\documentclass[conference,10pt]{IEEEtran}
\usepackage[utf8]{inputenc}
\usepackage{cite}
\usepackage{amsmath,amssymb,amsfonts}
\usepackage{algorithmic}
\usepackage{graphicx}
\usepackage{textcomp}
\usepackage{xcolor}
\usepackage{booktabs}
\usepackage{microtype}
\usepackage{hyperref}
\usepackage{listings}

\hypersetup{
    colorlinks=true,
    linkcolor=blue,
    citecolor=blue,
    urlcolor=blue
}

\begin{document}

\title{Evaluating Snapshot-Consistent Live Migration for Stateful Containers in Kubernetes: A CRIU and Sidecar-Free eBPF Study}

\author{
\IEEEauthorblockN{Musa Khasmammadov}
\IEEEauthorblockA{Independent Systems Researcher \\
Baku, Azerbaijan \\
musaxasmammedov77@gmail.com}
}

\maketitle

\begin{abstract}
Stateless microservices in Kubernetes tolerate node maintenance and cloud instance preemption through multi-replica rolling updates and PodDisruptionBudgets. Conversely, stateful singleton workloads that lack built-in distributed checkpointing---such as legacy numerical simulations, in-memory session aggregators, and warm machine-learning (ML) embedding caches---cannot be trivially replicated without consensus overhead or intrusive source-code refactoring. When an Amazon Web Services (AWS) EC2 Spot instance receives a two-minute interruption notice or a node is drained, the standard Kubernetes eviction sequence (\texttt{SIGTERM} followed by \texttt{SIGKILL}) destroys all volatile heap state.

We design, implement, and empirically evaluate an automated, snapshot-consistent live migration pipeline for stateful Kubernetes pods. Our system coordinates CRIU (v4.3) via the Kubelet Checkpoint API (KEP-2008), Buildah OCI layer packaging, the CRI-O runtime (v1.31), and a sidecar-free Cilium eBPF datapath. We evaluate the architecture on a three-node AWS EC2 cluster across four experimental dimensions:
(1) Across $N = 36$ live migrations of a 26~MB stateful service compared against $N = 20$ native cold reschedules, decoupled checkpoint migration achieves a median cutover blackout of 1064.5~ms (mean $1132.4 \pm 241.1$~ms, minimum 797~ms) versus 1019.0~ms (mean $1007.0 \pm 231.2$~ms) for cold restart---adding only 45.5~ms ($+4.5\%$) to median cutover latency and 144.0~ms of source freeze time while preserving in-memory state with a bounded Recovery Point Objective (RPO) of $669.0$~ms ($7.5 \pm 0.6$ ticks at 10~Hz).
(2) Linear regression across all 36 runs ($r = 0.9624$) demonstrates that CRI-O unfreezes the restored container in $\approx 537$~ms, whereas the remaining blackout variance arises from Kubelet's 1~s Pod Lifecycle Event Generator (PLEG) sync cycle.
(3) Scaling an in-memory numerical weight buffer from 26~MB to 1024~MB (1~GB) with cryptographic SHA-256 heap verification confirms 100\% bit-level state integrity across nodes within 50.09~s of source-side evacuation time ($41.7\%$ of the 120~s Spot deadline), while uncovering two systems bottlenecks: a Linux dirty-page writeback cliff above 512~MB and single-threaded \texttt{gzip} compression in Buildah ($32.06$~s at 1~GB).
(4) An empirical comparison against a two-container loopback sidecar pod demonstrates why single-container KEP-2008 checkpointing breaks in-pod sidecar proxies but succeeds under host-level Cilium eBPF routing.
\end{abstract}

\begin{IEEEkeywords}
Kubernetes, Container Live Migration, CRIU, CRI-O, Cilium eBPF, Stateful Workloads, Spot Instances, Recovery Point Objective.
\end{IEEEkeywords}

\section{Introduction}
Kubernetes schedules and heals workloads under the assumption that pod instances are stateless and interchangeable \cite{burns2016borg}. When a worker node undergoes kernel maintenance (\texttt{kubectl drain}) or an ephemeral cloud virtual machine is reclaimed, stateless deployments maintain availability by launching replacement replicas before terminating old pods (\texttt{maxSurge}).

However, many production workloads maintain large, mutable volatile heaps without native distributed persistence:
\begin{enumerate}
    \item \textbf{Single-Process Numerical and Scientific Runtimes:} Iterative solvers, physics engines, and legacy C++/Python simulations that accumulate state in RAM over hours of execution.
    \item \textbf{Warm ML Inference Caches and Session Aggregators:} In-memory embedding tables, KV caches, and stateful stream aggregators where rebuilding memory structures from cold object storage incurs severe warm-up penalties.
\end{enumerate}

Running such workloads on Amazon Web Services (AWS) EC2 Spot instances reduces compute expenditure by 70--90\% compared to On-Demand instances, but exposes pods to preemption with a hard 120-second termination notice (often preceded by an early EC2 Rebalance Recommendation) \cite{netflix_spot, aws_spot}. During standard eviction, Kubelet sends \texttt{SIGTERM}, waits for \texttt{terminationGracePeriodSeconds}, and issues \texttt{SIGKILL}, wiping the process address space. Offloading state to an external key-value store such as Redis \cite{redis_latency} avoids data loss during node preemption, but requires refactoring internal application data structures and replaces nanosecond-scale RAM references with sub-millisecond network serialization.

Userspace checkpoint/restore via CRIU \cite{criu} provides an application-transparent alternative \cite{stoyanov2018, nadgowda2017}. While Kubernetes v1.25+ introduced the Kubelet Checkpoint API (KEP-2008) primarily for forensic inspection \cite{reber2022}, orchestrating KEP-2008 into a complete cross-node live migration pipeline exposes critical systems questions around downtime decomposition, state rollback, memory scaling, and service-mesh compatibility.

This paper makes four contributions:
\begin{itemize}
    \item \textbf{End-to-End Orchestration Pipeline:} We build a privileged DaemonSet operator (\texttt{PodLiveMigration} CRD) that automates KEP-2008 checkpointing, Buildah OCI packaging, registry transfer, CRI-O restoration, and Cilium eBPF routing across AWS EC2 nodes without application code changes.
    \item \textbf{Statistical \& Regression Analysis ($N = 36$ vs. $N = 20$):} Comparing 36 CRIU migrations against 20 cold reschedules, we show that cutover blackout is governed in both cases by Kubelet's 1~s PLEG relist interval ($\sigma \approx 235$~ms). Regression of post-restore counter ticks against restore duration ($r = 0.9624$) isolates true container unfreeze time at $\approx 537$~ms.
    \item \textbf{Empirical Heap Scaling to 1~GB with SHA-256 Verification:} We migrate dense numerical heaps from 26~MB to 1024~MB with 100\% SHA-256 integrity, identifying the Linux \texttt{vm.dirty\_background\_ratio} writeback threshold and Buildah's single-threaded \texttt{gzip} step as the primary scaling bottlenecks.
    \item \textbf{Empirical Sidecar vs. Sidecar-Free Evaluation:} We experimentally demonstrate how KEP-2008's single-container checkpoint scope and \texttt{tcp-close} semantics sever in-pod sidecar loopback connections upon target restoration, validating the necessity of host-level eBPF service routing.
\end{itemize}

\section{System Architecture}

\subsection{Migration Pipeline and CRIU Configuration}
Our testbed runs on Ubuntu 24.04 LTS (Linux 6.8) with CRIU v4.3 \cite{criu}, Kubelet v1.31.14 (KEP-2008 enabled) \cite{reber2022}, CRI-O v1.31.4 configured with the \texttt{runc} runtime, Buildah, and Cilium v1.16.3 (\texttt{kubeProxyReplacement=true}). Figure~\ref{fig:pipeline} illustrates the nine-step migration workflow implemented in our operator (\texttt{05-live-migration-operator.yaml}).

\begin{figure*}[t]
\centering
\fbox{\parbox{0.95\textwidth}{
\textbf{Nine-Stage Decoupled Container Live Migration Pipeline:} \\
1. \textit{Hazard Trigger ($T_{\text{detect}}$):} Operator detects node maintenance (\texttt{cordon}), Spot Interruption Notice, or manual CRD creation. \\
2. \textit{Pre-Dump Probe \& Checkpoint Freeze ($T_{\text{freeze}}$):} Operator reads pre-dump counter $c_{\text{dump}}$ and invokes \texttt{POST :10250/checkpoint/\{ns\}/\{pod\}/\{container\}}. CRIU halts threads via \texttt{ptrace(PTRACE\_SEIZE)}, writes memory pages to \texttt{checkpoint.tar}, and unfreezes the source container (\texttt{KeepRunning: true}). \\
3. \textit{OCI Layer Packaging ($T_{\text{build}}$):} Buildah packages \texttt{checkpoint.tar} into a compressed OCI image annotated with \texttt{io.kubernetes.cri-o.annotations.checkpoint.name}. \\
4. \textit{Registry Push ($T_{\text{push}}$):} Buildah pushes the image to the in-cluster registry (\texttt{10.10.1.10:5000}). \textbf{Source-side evacuation is now complete.} \\
5. \textit{Pre-Kill Probe \& Source Deletion ($t_{\text{kill}}$):} Operator records $c_{\text{kill}}$ and issues \texttt{DELETE /pods/\{pod\}?gracePeriodSeconds=0}; cutover blackout ($T_{\text{blackout}}$) begins. \\
6. \textit{Target Pod Creation:} Operator creates the target Pod manifest pinned to the destination worker via \texttt{spec.nodeName}. \\
7. \textit{CRI-O Pull \& Restore ($T_{\text{unfreeze}} \approx 537$~ms):} Target CRI-O pulls the checkpoint layer, verifies identical bind mounts (\texttt{automountServiceAccountToken: false}), and invokes \texttt{runc/criu restore}. \\
8. \textit{Process Resumption:} Target process re-binds \texttt{0.0.0.0:8090} in its new Pod network namespace and immediately resumes execution from the checkpointed state. \\
9. \textit{PLEG Sync \& Readiness Probe ($T_{\text{restore}}$):} Kubelet PLEG updates \texttt{Pod.status.phase = Running} in \texttt{kube-apiserver}; operator verifies \texttt{HTTP 200 OK} on the new Pod IP and records $c_{\text{restore}}$.
}}
\caption{Execution sequence of the decoupled snapshot-consistent live migration pipeline.}
\label{fig:pipeline}
\end{figure*}

Because a pod restored on a different worker node receives a new Pod IP from the target node's CIDR (\texttt{10.244.2.0/24} $\to$ \texttt{10.244.1.0/24}), \texttt{/etc/criu/default.conf} specifies five options \cite{criu}:
\begin{itemize}
    \item \texttt{tcp-close}: Restores listening sockets (\texttt{0.0.0.0:8090}) in the \texttt{LISTEN} state while closing established TCP sockets bound to the old Pod IP.
    \item \texttt{skip-in-flight}: Drops transient in-flight TCP handshakes occurring during 50~Hz probing at the instant of \texttt{ptrace} seizure.
    \item \texttt{skip-mnt /run/.containerenv}: Skips CRI-O's transient metadata mount so \texttt{CRImportCheckpoint} mount validation succeeds.
    \item \texttt{ext-unix-sk} and \texttt{shell-job}: Allows dumping external UNIX sockets and inherited stdio descriptors.
\end{itemize}

\subsection{Cutover Blackout Definition}
Our benchmark service uses \texttt{HTTP/1.0} (closing the TCP socket after each response) behind a Kubernetes Service ClusterIP and Cilium L7 Envoy Ingress without long-hold request buffering. Between source pod deletion ($t_{\text{kill}}$) and target endpoint readiness ($t_4$), incoming requests fail immediately with connection refused or \texttt{HTTP 503}. We define cutover blackout $T_{\text{blackout}}$ strictly as the unavailability window during which probes fail, measured both by the operator clock ($t_4 - t_{\text{kill}}$) and by an external 50~Hz ($\pm 20$~ms) prober as the maximum interval between consecutive \texttt{HTTP 200 OK} responses.

\section{Latency, State Drift, and Evacuation Model}

\subsection{Stage Timestamps and Disruption Decomposition}
The operator records monotonic timestamps around each pipeline step: $t_0$ (before checkpoint API call), $t_1$ (after checkpoint archive creation), $t_2$ (after Buildah commit), $t_3$ (after registry push), $t_{\text{kill}}$ (immediately after the pre-delete probe $c_{\text{kill}}$ and just before \texttt{DELETE /pods}), and $t_4$ (upon the first \texttt{HTTP 200 OK} from the restored pod after Kubelet reports \texttt{phase == Running}):
\begin{align}
T_{\text{freeze}} &= t_1 - t_0, \quad T_{\text{build}} = t_2 - t_1, \quad T_{\text{push}} = t_3 - t_2 \\
T_{\text{restore}} &= t_4 - t_3, \quad T_{\text{blackout}} = t_4 - t_{\text{kill}} = T_{\text{restore}} - t_{\text{probe}}
\end{align}
where $t_{\text{probe}} = t_{\text{kill}} - t_3 \approx 1$--$2$~ms is the latency of the pre-delete HTTP GET used to sample $c_{\text{kill}}$. Thus $T_{\text{blackout}}$ is deterministically $1$--$2$~ms shorter than $T_{\text{restore}}$.

Client-perceived disruption consists of the source-side freeze stutter $T_{\text{freeze}}$ (while threads are paused by \texttt{ptrace} and incoming packets wait in the kernel TCP backlog) plus the cutover blackout $T_{\text{blackout}}$:
\begin{align}
T_{\text{disruption}} &= T_{\text{freeze}} + T_{\text{blackout}} \\
T_{\text{total}} &= T_{\text{freeze}} + T_{\text{build}} + T_{\text{push}} + T_{\text{restore}}
\end{align}

\subsection{State Drift (RPO), Post-Restore Ticks, and Rollback}
In \textbf{Mode A (Decoupled, \texttt{KeepRunning: true})}, CRI-O unfreezes the source container at $t_1$, allowing it to serve traffic while Buildah packages and pushes the image. Any state mutations occurring between the snapshot ($S_{\text{ckpt}} \approx c_{\text{dump}} + 1$) and pod deletion ($t_{\text{kill}}$) are lost upon cutover. At mutation rate $\lambda$, the Recovery Point Objective (RPO) and lost update count $\Delta S_{\text{lost}}$ are:
\begin{align}
\text{RPO} &= T_{\text{build}} + T_{\text{push}} \\
\Delta S_{\text{lost}} &= c_{\text{kill}} - S_{\text{ckpt}} \approx \lambda \cdot (T_{\text{build}} + T_{\text{push}})
\label{eq:lost}
\end{align}
Because clients querying the source pod between $t_1$ and $t_{\text{kill}}$ observe state values up to $c_{\text{kill}}$ that revert to $S_{\text{ckpt}}$ after cutover, workloads with external side effects require idempotency keys or a write-ahead log (WAL).

In \textbf{Mode B (Synchronous Stop-and-Copy)}, terminating or halting the source container immediately at $t_1$ yields $\text{RPO} \equiv 0$, but extends service unavailability to the entire pipeline ($T_{\text{blackout}}^{\text{stop-and-copy}} = T_{\text{total}}$).

Crucially, on the target node, \texttt{runc restore} unfreezes the container at time $T_{\text{unfreeze}} < T_{\text{restore}}$, and the restored process immediately resumes incrementing from $S_{\text{ckpt}}$ at rate $\lambda$. However, the operator only probes $c_{\text{restore}}$ after Kubelet's PLEG loop (1~s relist period) updates \texttt{Pod.status.phase = Running} at $T_{\text{restore}} = T_{\text{unfreeze}} + \Delta t_{\text{pleg}}$. During $\Delta t_{\text{pleg}}$, the restored process accumulates $\Delta S_{\text{post-restore}}$ new ticks:
\begin{equation}
c_{\text{restore}} - c_{\text{dump}} \approx \lambda \cdot (T_{\text{restore}} - T_{\text{unfreeze}})
\label{eq:post_restore}
\end{equation}

\subsection{Source-Side Spot Evacuation Condition}
When AWS issues a 120-second Spot Interruption Notice, only the \textit{source node} is terminated at $t = 120$~s; target-side pull and restoration ($T_{\text{restore}}$) execute on the surviving destination node after $t_3$. Therefore, Spot survival requires only that \textbf{source-side evacuation} $T_{\text{evac}}(M)$ complete within 120~s:
\begin{align}
T_{\text{evac}}(M) = &\ T_{\text{detect}} + T_{\text{freeze}}(M) \nonumber \\
&+ T_{\text{build}}(M) + T_{\text{push}}(M) < 120\text{ s}
\label{eq:spot_evac}
\end{align}
where $T_{\text{detect}} \approx 2.5$~s for a 5~s IMDSv2 poll interval.

\section{Empirical Evaluation}

\subsection{Testbed Setup}
Experiments were conducted on a three-node AWS EC2 cluster in \texttt{eu-central-1}: one \texttt{c7i-flex.large} control-plane node (2~vCPU, 4~GB RAM, hosting the OCI registry at \texttt{10.10.1.10:5000} and Cilium Ingress) and two \texttt{m7i-flex.large} workers (2~vCPU, 8~GB RAM, EBS \texttt{gp3} storage). All 11 microservices of the Google Online Boutique suite ran continuously across the workers to generate background multi-tenant CPU, memory, and CNI contention.

\subsection{Experiment 1: Statistical Distribution ($N = 36$ CRIU vs. $N = 20$ Cold Baseline)}
Table~\ref{tab:stats} summarizes the empirical distribution across $N = 36$ CRIU migrations ($n = 20$ from \texttt{worker-1} to \texttt{worker-2}; $n = 16$ from \texttt{worker-2} to \texttt{worker-1}) and $N = 20$ native cold reschedules of the 26~MB stateful counter service ($\lambda = 10$~Hz). Table~\ref{tab:traces} compares two individual UI traces (\#35 and \#36) and Mode B against the cold baseline.

\begin{table*}[t]
\centering
\caption{Statistical Comparison: $N = 36$ CRIU Live Migrations vs. $N = 20$ Native Cold Reschedules (26~MB Service)}
\label{tab:stats}
\resizebox{0.98\textwidth}{!}{%
\begin{tabular}{@{}lccccccccc@{}}
\toprule
\textbf{Configuration / Subset} & \textbf{$T_{\text{freeze}}$ (ms)} & \textbf{$T_{\text{build}}$ (ms)} & \textbf{$T_{\text{push}}$ (ms)} & \textbf{RPO (ms)} & \textbf{$T_{\text{restore}}$ (ms)} & \textbf{$T_{\text{blackout}}$ (ms)} & \textbf{$T_{\text{disruption}}$ (ms)} & \textbf{$T_{\text{total}}$ (ms)} & \textbf{$c_{\text{kill}} - c_{\text{dump}}$ (ticks)} \\ \midrule
\textbf{CRIU \texttt{w1} $\to$ \texttt{w2} ($n=20$)} & & & & & & & & & \\
\quad Median ($p50$) / $p95$ & 144.0 / 158 & 368.5 / 400 & 293.0 / 310 & 664.0 / 703 & 1033.0 / 1533 & 1031.5 / 1532 & 1178.0 / 1673 & 1857.0 / 2358 & 7.5 / 8 \\
\quad Mean $\pm$ SD [Min--Max] & $143.3 \pm 9.4$ & $370.9 \pm 17.1$ & $293.6 \pm 12.5$ & $664.6 \pm 27.6$ & $1112.9 \pm 249.2$ & $1111.2 \pm 249.3$ [817--1687] & $1254.5 \pm 250.1$ [958--1836] & $1920.8 \pm 248.7$ & $7.6 \pm 0.6$ [7--9] \\ \midrule
\textbf{CRIU \texttt{w2} $\to$ \texttt{w1} ($n=16$)} & & & & & & & & & \\
\quad Median ($p50$) / $p95$ & 145.0 / 155 & 372.5 / 384 & 301.5 / 308 & 672.5 / 688 & 1092.0 / 1474 & 1090.5 / 1472 & 1229.0 / 1622 & 1887.0 / 2277 & 7.0 / 8 \\
\quad Mean $\pm$ SD [Min--Max] & $145.4 \pm 7.1$ & $366.6 \pm 16.4$ & $296.6 \pm 12.4$ & $663.2 \pm 27.9$ & $1160.2 \pm 235.9$ & $1158.8 \pm 235.7$ [797--1682] & $1304.1 \pm 238.1$ [944--1828] & $1968.8 \pm 238.0$ & $7.4 \pm 0.5$ [7--8] \\ \midrule
\textbf{CRIU Combined ($N=36$)} & & & & & & & & & \\
\quad Median ($p50$) / $p95$ & \textbf{144.0} / 156 & \textbf{370.5} / 389 & \textbf{298.5} / 310 & \textbf{669.0} / 696 & \textbf{1066.0} / 1533 & \textbf{1064.5} / 1532 & \textbf{1203.0} / 1673 & \textbf{1878.0} / 2358 & \textbf{7.0} / 8 (RAM Kept) \\
\quad Mean $\pm$ SD [Min--Max] & $144.2 \pm 8.4$ & $369.0 \pm 16.7$ & $294.9 \pm 12.4$ & $663.9 \pm 27.4$ & $1133.9 \pm 241.1$ & $1132.4 \pm 241.1$ [797--1687] & $1276.6 \pm 242.6$ [944--1836] & $1942.1 \pm 241.8$ & $7.5 \pm 0.6$ [7--9] \\ \midrule
\textbf{Native Cold Baseline ($N=20$)} & & & & & & & & & \\
\quad Median ($p50$) / $p95$ & N/A & N/A & N/A & Total Loss & 1019.0 / 1244 & \textbf{1019.0} / 1244 & \textbf{1019.0} / 1244 & \textbf{1019.0} / 1244 & \textbf{100\% State Lost} \\
\quad Mean $\pm$ SD [Min--Max] & N/A & N/A & N/A & Total Loss & $1007.0 \pm 231.2$ & $1007.0 \pm 231.2$ [712--1312] & $1007.0 \pm 231.2$ [712--1312] & $1007.0 \pm 231.2$ & Reset to 0 [1--6] \\ \bottomrule
\end{tabular}%
}
\end{table*}

\begin{table*}[t]
\centering
\caption{Detailed Trace Breakdown and Operating Mode Trade-offs (RTO vs. RPO)}
\label{tab:traces}
\resizebox{0.98\textwidth}{!}{%
\begin{tabular}{@{}lccccccccc@{}}
\toprule
\textbf{Mode / Trace ID} & \textbf{Path} & \textbf{$T_{\text{total}}$} & \textbf{$T_{\text{freeze}}$} & \textbf{$T_{\text{build}}$} & \textbf{$T_{\text{push}}$} & \textbf{$T_{\text{blackout}}$} & \textbf{$T_{\text{disruption}}$} & \textbf{RPO} & \textbf{$\text{Counter}(c_{\text{dump}} \to c_{\text{kill}} \to c_{\text{restore}})$} \\ \midrule
\textbf{Cold Reschedule (UI Trace)} & \texttt{w2} $\to$ \texttt{w1} & 1289~ms & N/A & N/A & N/A & 1289~ms & 1289~ms & Total Loss & $43544 \to 8$ (Lost 43536 ticks) \\
\textbf{CRIU Decoupled (Trace \#35)} & \texttt{w1} $\to$ \texttt{w2} & 1935~ms & 159~ms & 385~ms & 310~ms & 1078~ms & 1237~ms & 695~ms & $51 \to 60 \to 57$ ($\Delta S_{\text{lost}}\approx 8, \Delta S_{\text{post}}\approx 5$) \\
\textbf{CRIU Decoupled (Trace \#36)} & \texttt{w2} $\to$ \texttt{w1} & 1642~ms & 147~ms & 388~ms & 308~ms & 797~ms & 944~ms & 696~ms & $77 \to 85 \to 79$ ($\Delta S_{\text{lost}}\approx 7, \Delta S_{\text{post}}\approx 1$) \\
\textbf{Stop-and-Copy (Mode B, Med.)} & Bidirectional & 1878~ms & 144~ms & 370.5~ms & 298.5~ms & 1878~ms & 1878~ms & \textbf{0~ms} & Zero Drift ($T_{\text{blackout}} = T_{\text{total}}$) \\ \bottomrule
\end{tabular}%
}
\end{table*}

Three key findings emerge from Table~\ref{tab:stats} and Table~\ref{tab:traces}:
\begin{enumerate}
    \item \textbf{Negligible Cutover Overhead vs. Cold Startup:} Across $N = 36$ CRIU runs and $N = 20$ cold reschedules, median cutover blackout is $1064.5$~ms for CRIU versus $1019.0$~ms for cold startup---an overhead of only \textbf{+45.5~ms (+4.5\%)} to pull the $5.16$~MiB checkpoint layer and restore memory pages via CRIU, while preserving the process heap instead of resetting it to zero.
    \item \textbf{Kubelet PLEG Dominance and Unfreeze Regression:} Both CRIU restore ($\sigma = 241.1$~ms) and cold startup ($\sigma = 231.2$~ms) exhibit nearly identical standard deviations, matching the theoretical standard deviation of a uniform 1~s Kubelet PLEG relist cycle ($1000 / \sqrt{12} = 288.7$~ms). Ordinary least-squares regression of $(c_{\text{restore}} - c_{\text{dump}})$ against $T_{\text{restore}}$ across all 36 CRIU runs yields Pearson correlation $r = 0.9624$, slope $\hat{\lambda} = 9.96$~ticks/s (matching the 10.0~Hz counter), and x-intercept $T_{\text{unfreeze}} = 537.1$~ms. Thus, CRI-O actually unfreezes the restored container in $\approx 537$~ms; the remaining $261$--$1151$~ms of $T_{\text{restore}}$ is Kubelet PLEG status synchronization delay ($\Delta t_{\text{pleg}}$).
    \item \textbf{Resolution of Trace \#35 vs. Trace \#36:} Across all 36 runs, $c_{\text{kill}} - c_{\text{dump}}$ is tightly bounded to $7.5 \pm 0.6$ ticks (range $7$--$9$), confirming Eq.~(\ref{eq:lost}). In Trace \#35 ($51 \to 60 \to 57$, $T_{\text{restore}} = 1080$~ms), the restored container ticked $\approx 5$ times during its $543$~ms PLEG wait ($1080 - 537$~ms), whereas in Trace \#36 ($77 \to 85 \to 79$, $T_{\text{restore}} = 798$~ms), the $261$~ms PLEG wait produced only $\approx 1$--$2$ post-restore ticks.
\end{enumerate}

\subsection{Experiment 2: Empirical Heap Scaling (26~MB to 1024~MB) with SHA-256 Verification}
To evaluate realistic stateful workloads beyond a small counter, we implemented an in-memory numerical state workload that populates a contiguous heap buffer of $M \in \{64, 128, 256, 512, 1024\}$~MB with high-entropy pseudo-random \texttt{float64} weight blocks (simulating an in-memory ML embedding cache or scientific grid). Before checkpointing on \texttt{k8s-worker-1}, the workload computes a SHA-256 digest of the entire heap buffer; after restoration on \texttt{k8s-worker-2}, the \texttt{/live/verify} endpoint recomputes the SHA-256 digest over the restored memory. Table~\ref{tab:scaling} reports the empirical measurements.

\begin{table*}[t]
\centering
\caption{Empirical Memory Scaling on AWS EC2 (26~MB to 1024~MB Heap) with Cross-Node SHA-256 Heap Integrity Verification}
\label{tab:scaling}
\resizebox{0.98\textwidth}{!}{%
\begin{tabular}{@{}lccccccccc@{}}
\toprule
\textbf{Allocated Heap} & \textbf{OCI Layer Size} & \textbf{$T_{\text{freeze}}$ (Dump)} & \textbf{$T_{\text{build}}$ (Buildah)} & \textbf{$T_{\text{push}}$ (Registry)} & \textbf{RPO ($\Delta t_{\text{drift}}$)} & \textbf{Source $T_{\text{evac}}$ (\% 120s SLA)} & \textbf{$T_{\text{restore}}$} & \textbf{$T_{\text{total}}$} & \textbf{Restored SHA-256 Integrity} \\ \midrule
\textbf{26~MB} ($N=36$ Med.) & 5.16~MiB & 0.144~s & 0.371~s & 0.299~s & 0.669~s (7 ticks) & \textbf{0.81~s} (0.7\%) & 1.066~s & 1.878~s & Verified (\texttt{Counter} state) \\
\textbf{64~MB} (Weights) & 69.11~MiB & 0.216~s & 0.761~s & 0.615~s & 1.376~s (14 ticks) & \textbf{1.59~s} (1.3\%) & 0.951~s & 2.543~s & \textbf{100\% Match} (\texttt{7785bfc701411dcb}) \\
\textbf{128~MB} (Weights) & 133.11~MiB & 0.258~s & 1.082~s & 0.929~s & 2.011~s (21 ticks) & \textbf{2.27~s} (1.9\%) & 1.695~s & 3.964~s & \textbf{100\% Match} (\texttt{cce52f481cd5ca56}) \\
\textbf{256~MB} (Weights) & 261.13~MiB & 0.375~s & 5.796~s & 1.685~s & 7.481~s (75 ticks) & \textbf{7.86~s} (6.6\%) & 2.108~s & 9.964~s & \textbf{100\% Match} (\texttt{4b1b1359a0222290}) \\
\textbf{512~MB} (Weights) & 517.15~MiB & 0.760~s & 16.120~s & 3.267~s & 19.387~s (195 ticks) & \textbf{20.15~s} (16.8\%) & 7.807~s & 27.954~s & \textbf{100\% Match} (\texttt{ab028a06dcdc25bf}) \\
\textbf{1024~MB} (1~GB) & 1029.21~MiB & \textbf{7.464~s} & \textbf{32.060~s} & 10.570~s & 42.630~s (428 ticks) & \textbf{50.09~s} (41.7\%) & 17.328~s & 67.422~s & \textbf{100\% Match} (\texttt{c0da7edfafcf15a6}) \\ \bottomrule
\end{tabular}%
}
\end{table*}

Table~\ref{tab:scaling} reveals three important systems properties:
\begin{enumerate}
    \item \textbf{Bit-Exact Restoration up to 1~GB within Spot SLA:} All heap sizes up to 1024~MB restored with 100\% SHA-256 verification. At 1~GB, source-side evacuation ($T_{\text{freeze}} + T_{\text{build}} + T_{\text{push}}$) completed in \textbf{50.09~s} ($41.7\%$ of the 120~s AWS Spot budget), confirming that 1--2~GB containers can be evacuated before Spot reclamation.
    \item \textbf{The Kernel Dirty-Page Cache Cliff in $T_{\text{freeze}}$:} From 64~MB to 512~MB, $T_{\text{freeze}}$ remains sub-second ($216$~ms $\to$ $760$~ms, effective write rate $\approx 680$~MiB/s) because Kubelet writes \texttt{checkpoint.tar} into the Linux page cache. Between 512~MB and 1024~MB, the checkpoint crosses the kernel's \texttt{vm.dirty\_background\_ratio} threshold ($\approx 780$~MB on an 8~GB node), forcing synchronous writeback to EBS \texttt{gp3} ($\approx 138$~MiB/s) and increasing $T_{\text{freeze}}$ by $9.8\times$ to $7.464$~s.
    \item \textbf{Buildah Single-Threaded \texttt{gzip} Bottleneck:} At 1~GB, $T_{\text{build}}$ consumes $32.060$~s ($64.0\%$ of source evacuation time and $75.2\%$ of the $42.630$~s RPO drift, where $\Delta S_{\text{lost}} = 447 - 19 = 428$ ticks matches $\lambda \cdot \text{RPO} = 426.3$ ticks). Inspection confirms that \texttt{buildah commit} compresses layers using single-threaded \texttt{gzip} ($\approx 32.1$~MiB/s). Disabling compression or using multi-threaded \texttt{zstd} for cluster-local registries is a high-leverage optimization.
\end{enumerate}

\subsection{Experiment 3: In-Pod Sidecar Proxy vs. Sidecar-Free Cilium eBPF}
To test service-mesh compatibility empirically, we deployed a two-container pod (\texttt{bench-sidecar-pod}) on \texttt{k8s-worker-1} comprising \texttt{stream-server} (listening on \texttt{127.0.0.1:8090}) and an \texttt{envoy-sidecar} proxy (listening on \texttt{0.0.0.0:15006} and forwarding over a persistent loopback TCP socket to \texttt{127.0.0.1:8090}). Prior to checkpointing, probes to \texttt{:15006} returned \texttt{HTTP 200 OK} (\texttt{ram\_counter = 505}).

When triggered via \texttt{PodLiveMigration}, KEP-2008 dumped \texttt{stream-server} in $118$~ms ($T_{\text{build}} = 266$~ms, $T_{\text{push}} = 180$~ms), and CRI-O restored the container on \texttt{k8s-worker-2}. However, because KEP-2008 checkpoints only a single container per invocation and \texttt{tcp-close} closes the established \texttt{127.0.0.1:8090} loopback socket upon restoration, the restored backend was permanently isolated in the target pod namespace ($T_{\text{blackout}} = \infty$, timing out after 120~s). By contrast, Cilium's sidecar-free eBPF datapath maintains L4/L7 routing in the host kernel outside the container namespace, allowing all 36 single-container migrations to succeed without sidecar re-peering.

\section{Related Work}
Clark et al. \cite{clark2005} established iterative pre-copy live migration for virtual machines. For Linux containers, Stoyanov and Kollingbaum \cite{stoyanov2018} reduced CRIU filesystem transfer overhead using an image cache proxy, and Nadgowda et al. \cite{nadgowda2017} developed Voyager to coordinate CRIU memory dumps with union storage federation. Reber \cite{reber2022} introduced forensic container checkpointing into Kubernetes via KEP-2008. Our work extends KEP-2008 into an automated cross-node migration controller on CRI-O and Cilium eBPF \cite{cilium}, providing empirical characterization of Kubelet PLEG latency, page-cache writeback cliffs, Buildah compression bottlenecks, and sidecar incompatibility.

\section{Limitations and Future Work}
First, our workload uses \texttt{HTTP/1.0} with CRIU's \texttt{tcp-close} option; migrating persistent \texttt{HTTP/1.1} keep-alive, WebSocket, or UDP flows requires L7 ingress retry buffering or CRIU \texttt{--tcp-established} with static Pod IP routing. Second, for heaps $\ge 256$~MB, single-threaded \texttt{gzip} in Buildah inflates RPO ($7.48$--$42.63$~s); replacing OCI tar/gzip packaging with direct \texttt{criu page-server} memory streaming and iterative pre-dumping (\texttt{criu dump --pre-dump}) is the natural next step to achieve sub-second freeze and RPO on multi-gigabyte workloads.

\section{Conclusion}
We presented an empirical study of snapshot-consistent container live migration in Kubernetes using CRI-O, CRIU v4.3, Buildah, and sidecar-free Cilium eBPF. Across $N = 36$ migrations and $N = 20$ cold reschedules on AWS EC2, CRIU restoration adds only $45.5$~ms ($+4.5\%$) to median cutover blackout ($1064.5$~ms vs. $1019.0$~ms) while preserving 100\% of snapshot heap state across footprints up to 1~GB. Regression and scaling analyses identify Kubelet's 1~s PLEG loop, the Linux dirty-page writeback threshold, and Buildah \texttt{gzip} compression as the primary targets for future runtime optimization.

\section*{Artifact Availability}
All manifests, Ansible playbooks, operator code, and benchmark harnesses are available at: \url{https://github.com/musaxasmammedov7/final-research}.

\bibliographystyle{IEEEtran}
\begin{thebibliography}{00}
\bibitem{burns2016borg} B. Burns, B. Grant, D. Oppenheimer, E. Brewer, and J. Wilkes, ``Borg, Omega, and Kubernetes,'' \textit{Communications of the ACM}, vol. 59, no. 5, pp. 50--57, 2016.
\bibitem{netflix_spot} Netflix Technology Blog, ``Creating Your Own EC2 Spot Market,'' Netflix Inc., 2015.
\bibitem{aws_spot} Amazon Web Services, ``Spot Instance Interruptions and EC2 Rebalance Recommendations,'' AWS Documentation, 2024.
\bibitem{redis_latency} Redis Ltd., ``Redis Latency Problems Troubleshooting,'' \url{https://redis.io/docs/management/optimization/latency/}, 2024.
\bibitem{criu} Checkpoint/Restore In Userspace, ``CRIU CLI Options and Architecture,'' \url{https://criu.org}, 2024.
\bibitem{stoyanov2018} R. Stoyanov and M. J. Kollingbaum, ``Efficient Live Migration of Linux Containers,'' in \textit{High Performance Computing: ISC High Performance 2018 International Workshops}, LNCS, vol. 11203, Springer, Cham, 2018, pp. 184--193.
\bibitem{nadgowda2017} S. Nadgowda, S. Suneja, N. Bila, and C. Isci, ``Voyager: Complete Container State Migration,'' in \textit{Proc. IEEE 37th Int. Conf. on Distributed Computing Systems (ICDCS)}, 2017, pp. 2137--2142.
\bibitem{reber2022} A. Reber, ``Forensic Container Checkpointing in Kubernetes,'' \textit{Kubernetes Official Blog / CNCF}, Dec. 2022. [Online]. Available: https://kubernetes.io/blog/2022/12/05/forensic-container-checkpointing-alpha/
\bibitem{cilium} Cilium Authors, ``Cilium: eBPF-based Networking, Observability, and Service Mesh,'' Linux Foundation, 2024.
\bibitem{clark2005} C. Clark, K. Fraser, S. Hand, J. G. Hansen, E. Jul, C. Limpach, I. Pratt, and A. Warfield, ``Live Migration of Virtual Machines,'' in \textit{Proc. 2nd USENIX Symp. on Networked Systems Design and Implementation (NSDI)}, 2005, pp. 273--286.
\end{thebibliography}

\end{document}
